\documentclass[11pt,letterpaper,reqno]{amsart}
% BEGIN INLINE manuscript-style
\makeatletter
\NeedsTeXFormat{LaTeX2e}[2020/10/01]
% Use with \documentclass[11pt,letterpaper,reqno]{amsart}.
% This file controls presentation only. It supplies no mathematical content.

\RequirePackage[T1]{fontenc}
\RequirePackage{lmodern}
\RequirePackage{microtype}
\RequirePackage[letterpaper,margin=1in,headheight=15pt]{geometry}
\RequirePackage{amsmath,amsthm,amssymb,mathtools}
\RequirePackage{enumitem}
\RequirePackage{perpage}
\RequirePackage{longtable}
\RequirePackage{xcolor}
\definecolor{manuscriptred}{rgb}{0.6,0.1,0.1}
\definecolor{manuscriptblue}{rgb}{0.1,0.3,0.5}
\definecolor{manuscriptmagenta}{HTML}{60445C}
\RequirePackage[breaklinks=true,colorlinks=true,
  linkcolor=manuscriptred,citecolor=manuscriptblue,urlcolor=manuscriptmagenta]{hyperref}
\RequirePackage{cleveref}
\hypersetup{pdfauthor={}}

% Retain the AMS paragraph, title, heading and running-head conventions.
% Equation counters restart each section: 2.i, 2.ii, ... .
% Appendix sections keep their letters: A.i, A.ii, ... .
\numberwithin{equation}{section}
\renewcommand{\theequation}{\thesection.\roman{equation}}
% Prose references include "eq." and omit parentheses around the number.
\renewcommand{\eqref}[1]{\textup{eq.~\ref{#1}}}
\crefname{equation}{eq.}{eqs.}
\Crefname{equation}{eq.}{eqs.}
\creflabelformat{equation}{#2#1#3}
\setcounter{tocdepth}{1}
\renewcommand{\thefigure}{\Roman{figure}}
\renewcommand{\thetable}{\Roman{table}}
\MakePerPage{footnote}
\renewcommand{\thefootnote}{%
  \ifnum\value{footnote}>9\arabic{footnote}\else\fnsymbol{footnote}\fi}

% Upright statements. Main results are A, B, C; supporting statements
% share a section-local sequence. These end marks indicate typography,
% never proof completeness, validation, or a deferred-proof status.
\theoremstyle{definition}
\newtheorem{ms@theorem}{Theorem}
\renewcommand{\thems@theorem}{\Alph{ms@theorem}}
\newtheorem*{ms@theoremstar}{Theorem}
\newtheorem*{ms@theoremzero}{Theorem 0}
\newtheorem{ms@proposition}{Proposition}[section]
\newtheorem{ms@lemma}[ms@proposition]{Lemma}
\newtheorem{ms@definition}[ms@proposition]{Definition}
\newtheorem{ms@corollary}[ms@proposition]{Corollary}
\newtheorem*{ms@corollarystar}{Corollary}

\newcommand{\ms@majorstart}{%
  \pushQED{\qed}\renewcommand{\qedsymbol}{$\blacksquare\square$}}
\newcommand{\ms@minorstart}{%
  \pushQED{\qed}\renewcommand{\qedsymbol}{$\blacksquare$}}
\newenvironment{theorem}[1][]{\ms@majorstart\ms@theorem[#1]}
  {\popQED\endms@theorem}
\newenvironment{theorem*}[1][]{\ms@majorstart\ms@theoremstar[#1]}
  {\popQED\endms@theoremstar}
\newenvironment{theoremzero}[1][]{\ms@majorstart\ms@theoremzero[#1]}
  {\popQED\endms@theoremzero}
\newenvironment{proposition}[1][]{\ms@minorstart\ms@proposition[#1]}
  {\popQED\endms@proposition}
\newenvironment{lemma}[1][]{\ms@minorstart\ms@lemma[#1]}
  {\popQED\endms@lemma}
\newenvironment{definition}[1][]{\ms@minorstart\ms@definition[#1]}
  {\popQED\endms@definition}
\newenvironment{corollary}[1][]{\ms@majorstart\ms@corollary[#1]}
  {\popQED\endms@corollary}
\newenvironment{corollary*}[1][]{\ms@majorstart\ms@corollarystar[#1]}
  {\popQED\endms@corollarystar}
% The source's lemma/proposition variants are deliberately not carried over:
% choose one marker for each statement type rather than encode proof status.

\theoremstyle{remark}
\newtheorem{remark}[ms@proposition]{Remark}
\newtheorem*{remark*}{Remark}
% The ordinary amsthm proof environment retains its open square.

\crefname{ms@theorem}{theorem}{theorems}
\Crefname{ms@theorem}{Theorem}{Theorems}
\crefname{ms@proposition}{proposition}{propositions}
\Crefname{ms@proposition}{Proposition}{Propositions}
\crefname{ms@lemma}{lemma}{lemmas}
\Crefname{ms@lemma}{Lemma}{Lemmas}
\crefname{ms@definition}{definition}{definitions}
\Crefname{ms@definition}{Definition}{Definitions}
\crefname{ms@corollary}{corollary}{corollaries}
\Crefname{ms@corollary}{Corollary}{Corollaries}
\crefname{remark}{remark}{remarks}
\Crefname{remark}{Remark}{Remarks}

% Bibliography configuration for this manuscript.
\makeatother
% END INLINE manuscript-style

\title[Periodic Wilson states and reflection positivity]{Periodic Wilson states and reflection positivity}
\date{9 November 2026}
\begin{document}
\raggedbottom
\begin{abstract}
We study finite-period Gibbs states of the four-dimensional Wilson
\(\operatorname{SU}(2)\) model at fixed lattice spacing. A transfer-kernel
comparison and an exact conditional integration give a common period
for these states. Reflected local perturbations then compare their
expectations with those of reflection-positive states. Compact convex
separation yields reflection positivity for the entire periodic class,
including the coordinate and diagonal lattice mirrors and bounded
Borel half-space observables. We also derive domination by finite
chessboard seminorms.
\end{abstract}
\maketitle
\tableofcontents

\section{Introduction}

Finite Wilson measures admit reflection inequalities on suitable
periodic lattices. Passing such an inequality to a subsequential
infinite-volume limit preserves it for that limit. A periodic Gibbs
state is instead defined by its local conditional distributions and
a finite-index group of translation periods. These conditions do not
identify it with a chosen limit of finite periodic measures. The
question is whether the reflection inequalities hold throughout this
periodic class.

The proofs below separate the symmetry of the states from the
positivity of their reflection forms. A contraction of the transfer
kernel first controls an antisymmetric observable on two adjacent
layers. Exact integration against the Gibbs conditional density
extends this comparison to observables on arbitrarily many layers.
The resulting bond-reflection symmetry supplies a period common to
all periodic states. With this period fixed, folding a repeated local
source produces a reflection-positive comparison state. A separating
local observable would contradict this comparison if a periodic
state failed the reflection inequality.

The finite-volume action has the compact-link form used in
Ba{\l}aban's lattice construction \cite[pp.~251--252]{Bal87}.
The state comparison is proved here from the conditional kernels,
the positive convolution spectrum and finite reflection forms.
No assertion of uniqueness of the Gibbs state is needed. The
arguments hold at fixed spacing and fixed coupling. Their boundary
estimates may depend on the spacing.
They provide no continuum limit or mass-gap conclusion.

Constants denoted by \(C\) may change between estimates. Their
dependence is restricted to the fixed coupling, the lattice
geometry and the specified finite supports unless additional
parameters are displayed. All Haar measures have mass one.


\section{Wilson measures and local limits}\label{sec:lp-measures}

Let \(a>0\) be the lattice spacing, let \(\beta\ge0\) be the
coupling, and let \(G=\operatorname{SU}(2)\).
We use lattice units in the proofs, so the lattice is \(\mathbb Z^4\)
and the physical position of \(x\) is \(ax\).
A link field assigns \(U_\mu(x)\in G\) to the positive link
\((x,x+e_\mu)\). The reversed link has value \(U_\mu(x)^{-1}\).
On a torus we retain the base and direction as the link index.
For \(\mu<\nu\) the plaquette \(p=(x;\mu,\nu)\) is oriented first
along \(e_\mu\) and then along \(e_\nu\). Let \(U_p\) be its ordered
boundary product and set
\[
 a_p(U)=1-\tfrac12\operatorname{tr}U_p,\qquad
 w_\beta(g)=\exp\{-\beta(1-\tfrac12\operatorname{tr}g)\}.
\]
The trace is real on \(G\), and \(0\le a_p\le2\).
If \(\Gamma\) is a full-rank period lattice, the Wilson measure on
\(\mathbb Z^4/\Gamma\) is
\begin{equation}\label{eq:lp-wilson}
 d\mu_\Gamma(U)=Z_\Gamma(0)^{-1}\prod_p w_\beta(U_p)\prod_b dU_b .
\end{equation}
The two products are indexed by plaquettes and positive links,
respectively, and \(dU_b\) is normalized Haar measure. The density is strictly positive. The
bound \(e^{-2\beta\#\{p\}}\le Z_\Gamma(0)\le1\) follows by integrating
the pointwise bounds on the weight against product Haar.
Multiplying numerator and denominator by the same factor
\(e^{\beta\,\#\{p\}}\) gives the weight
\(\exp(\frac\beta2\sum_p\operatorname{Re}\operatorname{tr}U_p)\).
This is the compact-link Wilson action used by Ba{\l}aban
\cite[pp.~251--252, (0.2)]{Bal87}.

Write \(\tau_v x=x+v\) for translation by \(v\in\mathbb Z^4\).
For a lattice isometry \(s\), write \(sF(U)=F(s^*U)\), where
\((s^*U)(x,y)=U(sx,sy)\) on the infinite lattice.
On a quotient, apply \(s\) to an oriented lattice lift of the link
and then project its base and direction, inverting the variable
when the direction reverses. This uses the indexed link even when
two quotient links have the same endpoints.
Every isometry that descends to the torus preserves \(\mu_\Gamma\).
It permutes the Haar variables, inverting those whose orientation
changes. It preserves the plaquette trace because a change of base
conjugates the holonomy and reversal inverts it.
Site gauges act by \(U^g(x,y)=g(x)U(x,y)g(y)^{-1}\).
The plaquette holonomy is conjugated at its base and Haar measure is
invariant under both endpoint multiplications.

On the infinite lattice, let \(\Omega=G^{\mathcal E_+}\) for the countable
set \(\mathcal E_+\) of positive links. A local function depends on
finitely many link variables. The product topology makes \(\Omega\)
compact and metrizable. Continuous local functions are uniformly dense
in \(C(\Omega)\). To see this, enumerate the links and freeze
coordinates beyond the first \(N\). This changes the metric
\(\sum_j 2^{-j}\min(1,\|U_j-V_j\|)\) by at most \(2^{-N}\).

For a finite set \(D\) of positive link indices, let \(\mathcal P(D)\)
be the plaquettes whose boundaries use a link of \(D\) in either
orientation. Let \(\partial D\) be the positive link indices used
by those boundaries and lying outside \(D\). For an exterior configuration \(\eta\), the Wilson
conditional density on \(G^D\) is
\begin{equation}\label{eq:lp-dlr}
 h_D(V\mid\eta)=
 \frac{\exp[-\beta\sum_{p\in\mathcal P(D)}a_p(V\eta)]}
 {\int_{G^D}\exp[-\beta\sum_{p\in\mathcal P(D)}a_p(W\eta)]\,dW}.
\end{equation}
Here \(V\eta\) uses \(V\) on \(D\) and \(\eta\) outside it.
Denote the resulting replacement kernel by \(\gamma_D\).
Its denominator lies between \(e^{-2\beta|\mathcal P(D)|}\) and \(1\).
The density is continuous in \((V,\eta|_{\partial D})\).
For every continuous local \(F\), \(\gamma_D F\) is continuous and local,
by uniform continuity on the finite compact coordinate space.
Uniform approximation then gives \(\gamma_D C(\Omega)\subset C(\Omega)\).

A Wilson DLR state is a probability \(\nu\) with
\(\nu\gamma_D=\nu\) for all finite \(D\). This is equivalent to
\(\gamma_DF\) being the conditional expectation of \(F\) given the
exterior links. Indeed, for bounded exterior-measurable \(J\) one has
\(\gamma_D(JF)=J\gamma_DF\), and the measure equation applies to \(JF\).
The converse follows by integrating the conditional identity.
A periodic state is invariant under a finite-index translation
subgroup of \(\mathbb Z^4\).

Every Wilson DLR state is invariant under site gauges.
For a gauge with finite support, take \(D\) to contain all incident
links. The exterior is fixed, and the Haar changes of variables on
\(D\) preserve every plaquette trace in \eqref{eq:lp-dlr}.
The DLR equation therefore gives invariance under that gauge.
For a general site gauge and a local function, its action depends only
on gauges at the endpoints of the finitely many links in its support.
A finite-support gauge with these values gives the same integral.
Continuous local functions determine the measures, so invariance
extends to bounded Borel tests.

\begin{lemma}[Local Gibbs limits]\label{lem:lp-compactness}
Fix \(a>0\) and \(\beta\ge0\). Let \(\Gamma_i\) be torus period lattices
whose shortest nonzero periods tend to infinity. The periodically
lifted Wilson laws have weak subsequential limits on \(\Omega\).
Each limit is a Wilson DLR state and retains the lattice isometries
that preserve the torus laws eventually. In particular, a cofinal
sequence of cubic tori yields a translation- and signed-permutation-
invariant Wilson DLR state. For a fixed finite-index period subgroup
\(Q\), the \(Q\)-invariant Wilson DLR class is compact and convex.

The local DLR passage also applies after adding a fixed real
periodic interaction of bounded Borel local terms, with a
uniformly finite range and finitely many term types per period cell.
Use tori compatible with that interaction. Along a convergent
subsequence, expectations converge for every bounded Borel local
test. If the interaction is multiplied by \(t\), any weak limit
of its DLR states as \(t\to0\) is Wilson DLR, with the same local
Borel convergence.
\end{lemma}

\begin{proof}
Probability measures on a compact metric space are sequentially
weakly compact. Fix \(D\) and continuous local \(F\). The window, its
interaction neighborhood and the support of \(F\) embed injectively on
all sufficiently large tori. Every torus interaction meeting
the window then has a unique lifted occurrence meeting \(D\),
counting the given term indices.
Integrating the torus variables of \(D\) first therefore gives exactly
\(\widetilde\mu_{\Gamma_i}(\gamma_DF)=\widetilde\mu_{\Gamma_i}(F)\),
where the tilde denotes periodic lift. The identity is local in
\(F\). The lifted law generally has other copies of the \(D\)-links
in its infinite exterior. The two integrands are continuous, so the
identity passes to a weak limit. Uniform approximation and the
contraction of \(\gamma_D\) extend it to \(C(\Omega)\).
Thus the Borel measures \(\nu\gamma_D\) and \(\nu\) agree.
An isometry acts continuously on \(\Omega\), so its invariance
passes to the limit. The same closed affine conditions establish
compactness and convexity for a fixed \(Q\).

For a Borel perturbation, use the same local conditional identity.
The conditional density on any fixed window \(D\) is bounded by
the exponential of the oscillation of the sum of interaction terms
meeting \(D\). This gives a uniform bound for the marginal density
relative to Haar on \(G^D\). The bound passes to weak limits by
testing against nonnegative continuous functions. For a bounded
Borel function on \(G^D\), approximation in Haar \(L^1\) by
continuous functions now gives convergence of the expectations.
The perturbed kernel preserves bounded Borel locality, so its
local conditional identities pass to the limit. Equality of the
resulting probability measures gives the full DLR equations.

If the unscaled perturbation terms meeting \(D\) have total supremum
bound \(C_D\), the ratio of the perturbed to Wilson conditional
density lies between \(e^{-2|t|C_D}\) and \(e^{2|t|C_D}\), uniformly
in the exterior. The conditional identities thus pass to the
Wilson kernel as \(t\to0\), and the marginal bound is uniform
for bounded \(t\).
\end{proof}

For a fixed finite bond window \(D\), the marginal density on \(D\)
of an unperturbed Wilson law on a sufficiently large torus and its lift is
\(\int h_D(V\mid\eta)\,d\widetilde\mu_\Gamma(\eta)\), by the local
conditional identity. The same formula holds in a Wilson limit.
As \(V\) ranges over \(G^D\), the continuous functions
\(h_D(V\mid\cdot)\) form a compact family in supremum norm.
Weak convergence and a finite uniform net give uniform convergence of
these marginal densities along a convergent subsequence. They all
have the lower bound \(e^{-2\beta|\mathcal P(D)|}\), so the
ratios converge uniformly to one. Consequently, along a fixed cofinal
family, the infimum of the uniform relative marginal error over its
torus limit points tends to zero. Otherwise a violating subsequence
would have a weak limit in the same class and satisfy the preceding
uniform convergence. These statements concern fixed \(a,\beta,D\).

\section{Finite reflection forms}\label{sec:lp-finite-rp}

The reflection with normal \(n\) at height \(c\) is
\[
 \theta_c x=x-\frac{2(n\cdot x-c)}{|n|^2}n,
 \qquad
 \Theta_cF=\overline{\theta_cF}.
\]
We take \(n=e_\mu\) or \(e_\mu\pm e_\nu\) with \(\mu<\nu\).
Preservation of \(\mathbb Z^4\) requires and is implied by
\(c\in\frac12\mathbb Z\) in the coordinate case and \(c\in\mathbb Z\)
in the diagonal case, by the displayed formula. There are four
coordinate and twelve diagonal orientations. A coordinate cut is a
site cut at integral \(c\) and a bond cut at half-integral nonintegral
\(c\). In physical units the height offsets are multiplied by \(a\).
On the infinite lattice a closed-half algebra consists of bounded
Borel functions of links with both endpoints in that half.

For \(q=|n|^2\) a split period lattice has the form
\begin{equation}\label{eq:lp-split}
 \Gamma=\Gamma_\parallel\oplus\mathbb Z(2zn/q),
 \qquad
 \Gamma_\parallel\subset\mathbb Z^4\cap n^\perp,
 \qquad \operatorname{rank}\Gamma_\parallel=3 .
\end{equation}
Here \(z\) is a positive integer and the normal period vector is
integral. The two cuts on the torus have heights \(c,c+z\) modulo
\(2z\). The reflection interchanges the two intervening strips and
fixes their common sites. At height \(c+z\) the displacement is
exactly one normal period. On a torus a half-space link must have
a lift lying wholly in the corresponding closed strip. This
geometric convention distinguishes the two links of a shortest
periodic circle even when they have the same projected endpoints.
For diagonal cuts we take \(z\ge3\) and require every nonzero
tangential period to have Euclidean length greater than \(2\).
A plaquette has diameter \(\sqrt2\), so its vertices remain distinct
under this quotient. The height ranges of the plaquettes crossing
the two cuts are disjoint modulo the normal height period.
Coordinate cuts allow all \(z\ge1\).

\begin{proposition}[Finite reflection positivity]\label{prop:lp-finite-rp}
Let \(\beta\ge0\). On a split torus as above the Wilson law is
reflection invariant and
\[
             \mu_\Gamma(\Theta_cF\,F)\ge0
\]
for every bounded complex Borel function on either closed half.
This includes both coordinate cut types on even coordinate tori.
If \(E\) is bounded, real and measurable on one closed half, the
normalized law with additional density \(e^{E+\theta_c E}\)
has the same two properties.
\end{proposition}

\begin{proof}
Write \(s(g)\in S^3\subset\mathbb R^4\) for the unit quaternion
coordinates of \(g\in G\). Then
\(\frac12\operatorname{tr}(gh^{-1})=s(g)\cdot s(h)\), and
\begin{equation}\label{eq:lp-quaternion-gram}
 w_\beta(gh^{-1})=e^{-\beta}\sum_{m\ge0}\frac{\beta^m}{m!}
 \sum_{i_1,\ldots,i_m=1}^4
 \prod_{j=1}^m s_{i_j}(g)\prod_{j=1}^m s_{i_j}(h).
\end{equation}
This is a positive Gram expansion in real bounded features. The
coefficients sum to \(e^{3\beta}\), so the expansion and its finite
products converge absolutely and uniformly.

At a coordinate site cut no plaquette meets both strict halves.
The plaquettes in the two halves occur in reflected pairs, and the
plaquettes lying on the cuts depend only on the shared links.
When the normal period is two, every plaquette with normal
edges still lies in one strip by its geometric lift, even
though its vertices lie on the cuts.
Conditional on the shared links, the reflection form is an absolute
square integrated against a positive measure.

At a diagonal cut use a signed permutation to take \(n=e_1+e_2\).
A link changes height by \(0\) or \(1\) in absolute value, so no
link joins strict opposite sides. The only crossing plaquettes
lie in the \((1,2)\) plane and have vertex heights \(c-1,c,c+1,c\)
in cyclic order. Orient the two paths \(P_+,P_-\) joining the
two height-\(c\) vertices in the same direction. The boundary
holonomy is \(P_+P_-^{-1}\) up to base change and reversal.
The reflection exchanges the two paths, and
\eqref{eq:lp-quaternion-gram} applies. The links within the cuts
are fixed pointwise and can be conditioned on.
All remaining plaquettes lie in one closed half or in the cuts.
Apply this argument at both cuts.

At a coordinate bond cut let \(X,Y\) be the parallel link arrays
in the two adjoining layers, identified by reflection, and let
\(T_x\) be the crossing link from the \(X\) layer to the \(Y\)
layer at \(x\). Reverse a crossing link when necessary to give it
this orientation, which preserves its Haar measure.
Each crossing plaquette has trace equal to that of
\(X_{xy}T_yY_{xy}^{-1}T_x^{-1}\).
For \((T\cdot Y)_{xy}=T_xY_{xy}T_y^{-1}\), its kernel before
integrating the crossing links is
\(C(X,T\cdot Y)\), where
\[
 C(X,Y)=\prod_{xy}w_\beta(X_{xy}Y_{xy}^{-1}).
\]
The function \(C\) is invariant under simultaneous layer gauges,
since \((S\cdot X)_{xy}(S\cdot Y)_{xy}^{-1}\)
is the conjugate of \(X_{xy}Y_{xy}^{-1}\) by \(S_x\).
Haar invariance therefore gives
\[
 \int C(X,T\cdot Y)\,dT
       =\iint C(S\cdot X,T\cdot Y)\,dS\,dT .
\]
The product Gram expansion of \(C\) now averages each feature in
its own argument. Thus the integrated crossing kernel is again
a positive Gram sum. The crossing links are independent Haar
variables before the crossing factors are integrated. They enter
neither half-space tests nor the one-side factors. The crossing
variables at the two cuts are distinct, so their kernels multiply.
For a coordinate period of two, the two layers may occur at both
cuts, but the crossing links at those cuts still have distinct
indices and the product Gram argument still applies.

In each case, after identifying the two halves by reflection,
the unnormalized pairing of \(F,H\) is a sum of terms
\(b_j\int\overline{I_j(F;Y_0)}I_j(H;Y_0)\,d\lambda(Y_0)\)
with \(b_j\ge0\). Here \(Y_0\) is the shared link array,
\(\lambda\) includes its plaquette factors, and \(I_j\) integrates
one half against its plaquette factors and the corresponding Gram
feature. There is no shared link array at bond cuts.
The coefficient bound and boundedness of the tests justify Fubini.
This proves positivity and Cauchy--Schwarz for the reflection pairing.
Reflection invariance follows from the lattice isometry in
\eqref{eq:lp-wilson}. For the perturbed law, insert \(Fe^E\) into
the pairing. Its additional density is reflection invariant.
\end{proof}

For a cofinal sequence of even cubic tori, every fixed continuous
local test in an infinite coordinate half-space eventually lies in
the corresponding torus half with the same link indices. Its reflection
inequality passes to any weak limit. More generally this passage
holds for each fixed mirror on cofinal split tori, and for the
reflected perturbations of Proposition~\ref{prop:lp-finite-rp}.
The following lemma completes the half-space algebra for these
limits and for the periodic states considered later.

\begin{lemma}[Completion of the half-space algebra]\label{lem:lp-borel-rp}
Let \(\nu\) be a probability measure on \(\Omega\) invariant under
an admissible reflection \(\theta_c\). If
\(\nu(\Theta_cF F)\ge0\) for every continuous local function of
the links in the closed positive half-space, the same inequality
holds for every bounded Borel function of those links and for
every bounded Borel function of the closed negative half-space.
\end{lemma}

\begin{proof}
Continuous functions are dense in \(L^2\) of the positive-half
marginal by regularity on its compact metric configuration space.
Continuous cylinder functions approximate them uniformly. Thus
each bounded Borel half-space function \(F\) has continuous local
approximants \(F_j\to F\) in \(L^2(\nu)\). Reflection invariance
makes \(\Theta_c\) an antilinear isometry. Cauchy--Schwarz gives
\[
 |\nu(\Theta_cF_j F_j)-\nu(\Theta_cF F)|
 \le \|F_j-F\|_2(\|F_j\|_2+\|F\|_2),
\]
so positivity passes to \(F\). A test on the negative half is
the reflected conjugate of a positive-half test. Its reflection
form has the same value, which proves the other assertion.
\end{proof}

\section{Chessboards and a pressure comparison}\label{sec:lp-chessboard}

\begin{lemma}[One-direction chessboard inequality]\label{lem:lp-chessboard}
In a split torus let the normal height period be \(2MH\), with
positive integers \(M\) and \(H\). Cut it into \(2M\) slabs
with faces \(c+kH\), each of an admissible type and with the
finite reflection forms of Proposition~\ref{prop:lp-finite-rp}.
Suppose \(f_k\) is a bounded complex Borel function of links
whose lattice lifts lie in closed slab \(k\). Copy a slab function
using the unique element of the group generated by the face
reflections that carries its slab to each destination. Conjugate
the function when that element reverses orientation.
Write \(d(f_k)\) for the Wilson expectation of this disseminated
product. Then \(d(f_k)\ge0\) and
\begin{equation}\label{eq:lp-chessboard}
 \left|\mu_\Gamma\left(\prod_{k=0}^{2M-1}f_k\right)\right|
                  \le \prod_{k=0}^{2M-1}d(f_k)^{1/(2M)} .
\end{equation}
\end{lemma}

\begin{proof}
Reflection in face \(j\) sends the slab index \(k\) to
\(2j-1-k\) modulo \(2M\). The generated group acts freely and
transitively on slabs. Its translations change indices by even
integers, and its reflections send an even index to an odd one. The corresponding
maps of the torus satisfy these relations, since \(2MH\) is its
normal height period. This gives unique copies from a reference
slab, with conjugation on the reflected copies. A constant
dissemination is the product of a function on one half with
its conjugated reflection, so it has nonnegative expectation.
Shared links may occur in two slab functions. Both halves permit
those links and the same identity applies.

Transport the \(2M\) profiles to the reference slab. For a filling
\(\sigma\) of the slabs by these profile labels, divide its expectation
by the product over slabs of \((d(f_{\sigma(k)})+\epsilon)^{1/(2M)}\),
where \(\epsilon>0\). Call the resulting number \(T(\sigma)\).
Reflection Cauchy--Schwarz gives
\[
     |T(\sigma)|^2\le T(\sigma_+)T(\sigma_-),
\]
where \(\sigma_\pm\) retain one semicircle each and reflect it
to fill the other. The denominator factors for the two
reflected fillings together repeat each original factor twice.
Both reflected values are nonnegative.
Choose a filling of maximal absolute value. If the maximum is
positive, both reflections of this filling again attain it.
Start at the first face of any constant cyclic run and retain
the half beginning there. If the run has length less than \(M\),
the retained half and its reflection have a run of twice that
length. Once the run has length at least \(M\), retain a full
semicircle of it to get a constant maximizing filling. A
constant normalized filling is at most one. Thus the maximum
is at most one, also when it is zero. Let \(\epsilon\downarrow0\).
\end{proof}

\begin{corollary}[Cubic chessboards]\label{cor:lp-board-chessboard}
Let \(\Gamma=\bigoplus_{\mu=1}^4N_\mu\mathbb Z e_\mu\), where
the positive integers \(N_\mu\) are the coordinate periods. Tile
\(\mathbb Z^4/\Gamma\) by congruent axis-aligned cells whose width
in direction \(\mu\) divides \(N_\mu\), giving an even number
\(n_\mu\) of cells in that direction. Place their faces at
coordinate bond cuts. A cell function depends only on links wholly within that cell.
Let \(N=\prod_\mu n_\mu\), and define the chessboard seminorm of
a cell function by
\[
 p_\Gamma(F)=\mu_\Gamma\!\left(\prod_C F^{[C]}\right)^{1/N},
\]
where the copies are disseminated in the four coordinate directions.
For bounded complex Borel cell functions \(F_C\),
\begin{equation}\label{eq:lp-board}
 \left|\mu_\Gamma\left(\prod_C F_C\right)\right|
        \le\prod_Cp_\Gamma(F_C).
\end{equation}
Moreover \(p_\Gamma(1)=1\), \(p_\Gamma(\bar F)=p_\Gamma(F)\),
and \(p_\Gamma\) is absolutely homogeneous and subadditive.
\end{corollary}

\begin{proof}
Apply Lemma~\ref{lem:lp-chessboard} first to the products in the
slabs of direction one. Each resulting expectation is the
dissemination of one such slab. Regroup its factors by their
slabs in direction two and apply the same lemma. Reflections
in distinct directions commute, and conjugating a product
conjugates its factors. Thus the second step copies the previous
pattern in the second direction, with the product of the two
root exponents. Repeat in directions three and four.
Each original cell then has exponent \(1/N\).
The fully disseminated expectation is nonnegative by splitting
at any pair of opposite faces. Exactly half the copies are
conjugated, giving absolute homogeneity. Conjugating the entire
nonnegative expectation gives the conjugation identity.
Expand the dissemination of \(F+G\). Applying \eqref{eq:lp-board}
to each term gives the binomial bound
\((p_\Gamma(F)+p_\Gamma(G))^N\). Taking the root proves subadditivity.
\end{proof}

\begin{lemma}[Gibbs lower tangent]\label{lem:lp-lower-tangent}
Let \(Q\subset\mathbb Z^4\) have finite index and let \(A\) be a
bounded real Borel local observable. Let \(\Gamma_i\subset Q\)
have rectangular fundamental parallelepipeds \(D_i\) with all
side lengths tending to infinity. Put
\(\kappa=[\mathbb Z^4:Q]^{-1}\) and \(V_i=[\mathbb Z^4:\Gamma_i]\).
Let \(Z_i(t)\) be the integral of the unnormalized Wilson weight
multiplied by
\(\exp(t\sum_{v\in Q/\Gamma_i}(\tau_vA)^{\Gamma_i})\), with the
superscript denoting periodization of the observable.
For every \(Q\)-invariant Wilson DLR state \(\nu\) and fixed real \(t\),
\begin{equation}\label{eq:lp-lower-tangent}
        V_i^{-1}\log\frac{Z_i(t)}{Z_i(0)}
                 \ge t\kappa\,\nu(A)-o(1).
\end{equation}
The error is uniform in \(\nu\). The statement applies to the
rectangular split tori in \eqref{eq:lp-split}, including diagonal
normals.
\end{lemma}

\begin{proof}
Take \(D_i\) half-open, so its lattice vertices represent the
torus classes once. Let \(\mathcal L_i\) consist of the positive
links based at those vertices. This identifies its Haar variables
with the torus link variables. Since \(\Gamma_i\subset Q\), the
points of \(Q\cap D_i\) represent \(Q/\Gamma_i\). Let \(W_i\)
sum over those finitely many anchors whose translated source
supports are strictly interior to the window. This convention
also includes constant sources, whose supports are empty.
With the exterior links \(\eta\) fixed, Jensen's inequality for
the conditional Wilson law gives
\[
 \log\gamma_{\mathcal L_i}(e^{tW_i})(\eta)
                        \ge t\gamma_{\mathcal L_i}W_i(\eta),
\]
The logarithm on the left is an open-boundary partition ratio.
Away from a collar of fixed thickness the torus and open
plaquette and source terms agree as functions of the common
Haar variables. In the collar their absolute differences have
sum at most \(C(\beta+|t|\|A\|_\infty)|\partial_R D_i|\), where
\(|\partial_R D_i|\) counts the sites in a collar of thickness
\(R\) on both sides of the boundary. Each plaquette term is
bounded by \(2\beta\) and each fixed translated support has
bounded incidence per link. Comparing the two numerators and
denominators by these pointwise bounds gives the same order of
error for the logarithmic ratios, uniformly in \(\eta\).

There are \(\kappa V_i\) torus source copies, and all but
\(O(|\partial_R D_i|)\) occur in \(W_i\).
Integrating the Jensen bound against \(\nu\), the DLR equation
and \(Q\)-invariance give
\[
 \nu(\gamma_{\mathcal L_i}W_i)=\nu(W_i)
       =(\kappa V_i+O(|\partial_R D_i|))\nu(A).
\]
Finally \(V_i\) is the Euclidean volume of \(D_i\).
Disjoint unit cubes centered at the collar lattice sites lie
in a collar enlarged by a fixed constant, whose volume divided
by \(V_i\) tends to zero as every side length tends to infinity.
This also proves the claim for tilted rectangular boxes.
\end{proof}

\section{Bond reflections and a common period}\label{lp:sec-link-symmetry}

The period of a periodic Gibbs state may initially depend on the state.
The next theorem replaces it by a common period. The proof uses a
perturbation of the finite transfer kernel to compare adjacent layers,
then uses the conditional law to remove outer layers of an arbitrary
finite observable.

\begin{theorem}[Bond-reflection symmetry]\label{lp:thm-link-symmetry}
At spacing \(a>0\) and coupling \(\beta>0\), every periodic Wilson
DLR state is invariant under reflection in every coordinate bond
mirror \(x_\mu=a(i+\tfrac12)\), where \(i\in\mathbb Z\). In particular it is
invariant under all translations in \(2a\mathbb Z^4\).
\end{theorem}

Assume \(\beta>0\) throughout this section and fix one coordinate
as time. Section~\ref{sec:lp-measures} proves
gauge invariance of every Wilson DLR state. A reflected state has
the same property, since reflection preserves the Wilson specification.

\subsection{The transfer kernel}

Let \(C\) be the convolution operator on \(L^2(G)\) with kernel
\(w_\beta(uv^{-1})\). This kernel is real, symmetric, and at least
\(e^{-2\beta}\). For each integer \(d\ge1\), let \(\Pi_d\) be the
Peter--Weyl projection for the irreducible representation of dimension
\(d\). Convolution by a central function acts by a scalar on each
such space. The Haar class density \(2\sin^2 t/\pi\) and the
character \(\sin(dt)/\sin t\) give
\begin{equation}\label{eq:lp-convolution-spectrum}
 C=\sum_{d\ge1}\lambda_d\Pi_d,\qquad
 d\lambda_d=e^{-\beta}(I_{d-1}(\beta)-I_{d+1}(\beta)),\qquad
 \lambda_d=\frac{2e^{-\beta}}{\beta}I_d(\beta)>0 .
\end{equation}
Here
\[
 I_n(\beta)=\frac1\pi\int_0^\pi e^{\beta\cos t}\cos(nt)\,dt
     =\sum_{k\ge0}\frac{(\beta/2)^{n+2k}}{k!(n+k)!}.
\]
The series follows by expanding
\(e^{(\beta/2)(e^{it}+e^{-it})}\). Integration by parts in
\(\int_0^\pi e^{\beta\cos t}\sin t\sin(dt)\,dt\) gives the
last equality in \eqref{eq:lp-convolution-spectrum}.
The series bounds \(I_{n+1}/I_n\) by \(\beta/(2(n+1))\), while
\[
 \frac{I_n}{I_{n+1}}
     =\frac{2(n+1)}{\beta}+\frac{I_{n+2}}{I_{n+1}}
     \le C_\beta(1+n).
\]
Thus the ratios between indices of bounded distance grow at most
polynomially in the indices. The factorial decay of \(\lambda_d\),
with the polynomial dimensions and derivatives of the matrix
coefficients, gives trace class operators \(C\) and \(C^{1/2}\) with
smooth real symmetric kernels.

For a finite set \(\mathcal B\) of positive spatial bonds use
\(C_{\mathcal B}=\bigotimes_{e\in\mathcal B}C\). Its kernel is
\[
 C_{\mathcal B}(u,v)=\prod_{e\in\mathcal B}w_\beta(u_ev_e^{-1}).
\]
Simultaneous vertex gauge transformations preserve this kernel. Write
\(\Pi_{\boldsymbol d}\) for the tensor projections and
\(\lambda_{\boldsymbol d}=\prod_e\lambda_{d_e}\).
For a polynomial \(f\) in matrix entries and their conjugates,
define on smooth functions
\begin{equation}\label{eq:lp-smooth-transfer}
 \widehat f F
 =\sum_{\boldsymbol d,\boldsymbol d'}
   \frac{\lambda_{\boldsymbol d}}{\lambda_{\boldsymbol d'}}
   \Pi_{\boldsymbol d'}M_f\Pi_{\boldsymbol d}F .
\end{equation}
The tensor-product rule for \(\operatorname{SU}(2)\) confines each
coordinate difference \(|d_e-d'_e|\) to a bound determined by
the degree of \(f\) in that bond. Thus each index has a bounded
number of partners, and each displayed ratio grows at most
polynomially. The coefficients of a smooth function decay faster than
any power of the joint index, by applying powers of the Casimir in
each group coordinate. The sup norm of a matrix-coefficient block
and of each fixed derivative is bounded by its \(L^2\) norm times
a polynomial in the indices. This follows from orthogonality, the
block dimension, and differentiation in the finite-dimensional
representations. The series therefore converges in every smooth norm
and defines a continuous map on smooth functions. It also acts on
smooth families in finitely many additional variables, because
derivatives in those variables commute with each term and satisfy
the same bounds.

Multiplying the series by \(C_{\mathcal B}\) gives
\(C_{\mathcal B}\widehat f=M_f C_{\mathcal B}\) on this domain.
By symmetry of the kernel this is the pointwise identity
\begin{equation}\label{eq:lp-kernel-transfer}
 f(w)\int C_{\mathcal B}(v,w)F(v)\,dv
     =\int C_{\mathcal B}(v,w)(\widehat fF)(v)\,dv .
\end{equation}
The spectral definition in \eqref{eq:lp-smooth-transfer} supplies
the domain of this operation. In particular the proof uses no
application of an unbounded convolution inverse to a general
smooth function.

Consider a cubic torus of integer side \(s\ge2\) in lattice units,
and let \(\mathcal E\) be the spatial bonds in one layer. Let \(P_{\rm g}\)
average over the vertex gauge group of that layer, and let \(M\)
multiply by its spatial Wilson factor. Define the transfer operator by
\[
 T=P_{\rm g} C_{\mathcal E}^{1/2}M C_{\mathcal E}^{1/2}P_{\rm g}.
\]
It is positive and trace class. The gauge action preserves the spatial
Wilson factor and commutes with the convolution, so both \(M\) and
\(C_{\mathcal E}^{1/2}\) commute with \(P_{\rm g}\). Across a time step the
two spatial arrays \(u,v\) and the time links \(g\) give the
factor \(C_{\mathcal E}(u,g\cdot v)\). On the bond \((x,y)\)
the corresponding time plaquette has product
\(u_{xy}g_yv_{xy}^{-1}g_x^{-1}=u_{xy}((g\cdot v)_{xy})^{-1}\).
Integrating \(g\) inserts
the projection \(P_{\rm g}\), with no extra factor because Haar is
normalized. The remaining cyclic kernel integral equals
\(\operatorname{Tr}(T^s)\).

Let \(\Xi\) be a bounded operator commuting with complex conjugation on
\(L^2(G^{\mathcal B})\) of norm at most one, commuting with the
vertex gauge action. Suppose \(C_{\mathcal B}^{1/2}\Xi
C_{\mathcal B}^{1/2}\) has a continuous strictly positive kernel
\(k_\Xi(u,v)\). Define the real local perturbation
\begin{equation}\label{eq:lp-transfer-perturbation}
 \phi_\Xi(u,v)=\log\frac{k_\Xi(u,v)}{C_{\mathcal B}(u,v)},
\end{equation}
where the two spatial arrays in the original field are compared
after transporting the upper one by the intervening time links.
The simultaneous gauge covariance makes this a gauge-invariant
observable of the original bonds.

Take a periodic DLR state \(\nu\) and choose \(m\mathbb Z^4\)
in its period group, with \(m\) large enough that translates
of the spatial block at a given time step are disjoint. Repeat
\(\phi_\Xi\) at those anchors on tori with \(s\) a sufficiently
large multiple of \(m\). At step \(j\), the modified crossing
operator before time-link integration is
\(C_{\mathcal E}^{1/2}\Xi_j C_{\mathcal E}^{1/2}\), where \(\Xi_j\)
is the tensor product of the inserted contractions and identities
on the other bonds. Each \(\Xi_j\) has norm at most one and
commutes with \(P_{\rm g}\). The modified partition function is the trace of the
cyclic product of \(M C_{\mathcal E}^{1/2}\Xi_j C_{\mathcal E}^{1/2}P_{\rm g}\).
The trace formulas follow by splitting the cyclic integral into
two Hilbert--Schmidt kernels. Cyclicity of the trace and
commutation with \(P_{\rm g}\) rewrite the product as \(\Xi_1T\cdots\Xi_sT\)
inside the trace. The Schatten H\"older inequality, with
\(\|\cdot\|_{\mathfrak S_s}\) the Schatten \(s\)-norm, gives
\begin{equation}\label{eq:lp-transfer-contraction}
 0<Z_s(\phi_\Xi)=\operatorname{Tr}(\Xi_1T\cdots\Xi_sT)
 \le\prod_{j=1}^s\|\Xi_jT\|_{\mathfrak S_s}
 \le\|T\|_{\mathfrak S_s}^s=Z_s(0).
\end{equation}
The positivity at the left comes from the kernels in the modified
integral, and the final equality uses positivity of \(T\).
Lemma~\ref{lem:lp-lower-tangent} with period \(m\mathbb Z^4\)
now gives \(\nu(\phi_\Xi)\le0\).

\subsection{The comparison of two adjacent layers}

Let \(Y(u,v)\) be real and continuous on the two arrays of a finite
spatial block \(\mathcal B\). Suppose it is invariant under a common
vertex gauge and antisymmetric under exchange of the arrays.
As in \eqref{eq:lp-transfer-perturbation}, transport the upper array
by the intervening time links. This defines the local observable of
the original bond field whose expectation is denoted by \(\nu(Y)\).
We claim
\begin{equation}\label{eq:lp-adjacent-antisymmetry}
                           \nu(Y)=0 .
\end{equation}
First suppose \(K(u,v)=Y(u,v)C_{\mathcal B}(u,v)\) is a polynomial.
It is the kernel of a finite-rank real skew-adjoint operator.
Choose a finite sum of joint Peter--Weyl spaces containing the
support in both kernel variables. On this space define
\(D=C_{\mathcal B}^{-1/2}K C_{\mathcal B}^{-1/2}\), using the
positive eigenvalues in \eqref{eq:lp-convolution-spectrum}, and
extend by zero on the orthogonal complement. The extension is real,
skew-adjoint and commutes with the vertex gauge action. Hence
\(\Xi_\varepsilon=e^{\varepsilon D}\) is an isometry for either
sign of the real parameter \(\varepsilon\). Its sandwiched kernel
is \(C_{\mathcal B}+\varepsilon K+O_K(\varepsilon^2)\) uniformly,
because the difference has finite spectral support. Strict
positivity of \(C_{\mathcal B}\) admits the kernel when
\(|\varepsilon|<\varepsilon_0(K)\). Here the remainder bound and
the positive radius depend on the fixed polynomial kernel \(K\),
as well as the fixed coupling and block. The logarithm in \eqref{eq:lp-transfer-perturbation}
therefore gives \(\phi_{\Xi_\varepsilon}=\varepsilon Y+O_K(\varepsilon^2)\)
uniformly. The pressure comparison at each fixed \(\varepsilon\)
gives \(\varepsilon\nu(Y)\le C_K\varepsilon^2\).
Let \(\varepsilon\) tend to zero through both signs.

For general \(Y\), approximate \(Y C_{\mathcal B}\) uniformly by
real matrix-coordinate polynomials. Average each approximant
under the common vertex gauge group and take its antisymmetric
part. Both operations preserve polynomiality and uniform convergence.
Division by the strictly positive kernel yields functions to which
the polynomial case applies, converging uniformly to \(Y\).
This proves \eqref{eq:lp-adjacent-antisymmetry}.

\subsection{Removing the outer layers}

On a finite stack of layers \([j_-,j_+]\), write \(t_{j,x}\) for
the time link from \((j,x)\) to \((j+1,x)\), and \(U_j\) for
the spatial array of layer \(j\) transported to the lowest layer.
Explicitly, at vertex \(x\) use
\(g_{j,x}=t_{j_-,x}\cdots t_{j-1,x}\) and act on the spatial
bonds by \(g_j\) at their endpoints. In these coordinates the
time links of the stack are the identity. A gauge-invariant
polynomial observable of the stack is therefore a polynomial of
the transported spatial arrays invariant under a common vertex
gauge. Only finitely many vertices are needed for any local test.

At an interior layer \(j\), resample the original spatial bonds
of a finite block \(\mathcal B\), conditioning on all other bonds.
The time links are fixed in this resampling, so the substitution
to the transported variables preserves Haar. With
\(w=U_{j-1}|_{\mathcal B}\), \(v=U_j|_{\mathcal B}\) and
\(w'=U_{j+1}|_{\mathcal B}\), the conditional density is proportional to
\begin{equation}\label{eq:lp-layer-density}
        C_{\mathcal B}(w,v)e^{-V(v;\eta)}
                         C_{\mathcal B}(v,w')\,dv .
\end{equation}
The two kernels account for the time plaquettes through these
spatial bonds. The potential \(V\) is \(\beta\) times the sum of
the spatial plaquette actions meeting the block, and \(\eta\)
denotes its same-layer exterior variables. This list exhausts
the plaquettes in the finite-bond conditional density.

For a polynomial \(f(w)\) and a smooth local factor \(\Phi\)
depending on the transported variables from layer \(j\) upward,
put
\begin{equation}\label{eq:lp-single-elimination}
 \Phi^\flat=
 \frac{\widehat f_v\bigl[e^{-V(v;\eta)}C_{\mathcal B}(v,w')\Phi\bigr]}
      {e^{-V(v;\eta)}C_{\mathcal B}(v,w')}.
\end{equation}
The denominator is positive on the compact variable space.
The smooth transfer operation leaves \(\Phi^\flat\) smooth and
confines its dependence to finitely many variables from layer
\(j\) upward. The new support can include same-layer neighbours and
the block in layer \(j+1\). The identity \eqref{eq:lp-kernel-transfer}
in the conditional numerator of \eqref{eq:lp-layer-density}
gives \(\nu(f(w)\Phi)=\nu(\Phi^\flat)\). The normalizing factor
is the same on both sides of that conditional equality. The
reflected construction removes the upper outer layer.

Consider a stack of \(2k+2\) layers symmetric about a bond mirror,
where \(k\ge1\). Express a real invariant polynomial of the
stack as a sum of terms \(f(w)h(z)\Phi\), using the outer arrays
\(w,z\) and a common finite spatial block. Write \(u,v\) for
the first and last retained arrays. Applying the lower then
the upper identity to a term gives the expectation of
\begin{equation}\label{eq:lp-double-elimination}
                 \Upsilon^{-1}\widehat f_u\widehat h_v
                                      (\Upsilon\Phi).
\end{equation}
When \(u,v\) are nonadjacent, \(\Upsilon\) is the product of
their spatial factors and their two kernels toward the interior.
For adjacent \(u,v\), the lower removal divides by
\(e^{-V(u;\eta)}C_{\mathcal B}(u,v)\). The upper conditional
numerator multiplies by \(e^{-V(v;\eta')}C_{\mathcal B}(u,v)\)
before \(\widehat h_v\) acts. The common kernel cancels in
these outside factors. The remaining multipliers independent
of the operator variables can be moved through the operators,
which commute on the distinct arrays. Thus in this case
\[
             \Upsilon=e^{-V(u;\eta)}e^{-V(v;\eta')}
                                   C_{\mathcal B}(u,v).
\]
This explains why the common kernel enters once. For a fixed block,
the expression is linear on the tensor product of the two outer
polynomial spaces and the inner smooth functions. These are spaces
of functions on independent variables, so the expression respects
any decomposition of the original polynomial into products.
Reversing
the order of the layers exchanges the two operator factors
and preserves the corresponding \(\Upsilon\), including its
neighbour variables. The two-sided elimination therefore
commutes with reversal of the transported array.

Average the output of \eqref{eq:lp-double-elimination} over a
common vertex gauge on its finite set of spatial sites. Its
expectation stays the same, since this action can be realized
by ordinary gauge transformations at the lowest original
layer. The averaged output no longer depends on the choice
of lowest layer for temporal transport. Indeed, changing
the base to the first retained layer applies one common
gauge to all retained arrays. Under the original bond
reflection the transported array is reversed and transformed
by the common gauge \(g_{j_+}^{-1}\). Both statements follow
by multiplying the temporal links in the corresponding
order. Gauge invariance of the averaged output removes
these common actions pointwise. We obtain a continuous
gauge-invariant observable of the retained stack, with
the same expectation as the original polynomial and with
elimination commuting with the original reflection.

\begin{proof}[Proof of Theorem~\ref{lp:thm-link-symmetry}]
For a gauge-invariant polynomial on the two layers adjacent to
the mirror, reflection exchanges the transported arrays up
to a common gauge. The difference of the two values is
antisymmetric and has zero expectation by
\eqref{eq:lp-adjacent-antisymmetry}. For an invariant polynomial
on \(2k+2\) layers, the two-sided identity and its reflected
version reduce the expectation comparison to a continuous
invariant function on \(2k\) layers. Matrix-coordinate
polynomials are uniformly dense among continuous functions
on the finite product of groups by Stone--Weierstrass.
Averaging the approximants over the relevant vertex gauge
group gives uniform density within the invariant functions.
Induction on \(k\) therefore proves reflection invariance of
the expectations of invariant continuous local tests.

Every finite local support fits in a finite even stack symmetric
about the fixed bond mirror. For a general continuous local test,
average over the vertex gauges of its finite support. The state and its reflected
state have the same expectations before and after this
averaging. The averaged test is invariant, so the two
states agree on continuous local tests. These determine the
measure on the product space. The time direction and the
bond mirror were arbitrary. Two consecutive parallel bond
reflections compose to translation by two lattice steps,
which gives the claimed common period in physical units.
\end{proof}

\section{Reflection positivity of periodic Gibbs states}
\label{lp:sec-fold}

Work in lattice units. By Theorem~\ref{lp:thm-link-symmetry},
the periodic Wilson DLR states form the class \(\mathcal G_\beta^P\)
of Wilson DLR states invariant under \(P=2\mathbb Z^4\). A folded
source will produce states in this class that dominate the
expectation of a prescribed local test.

\begin{lemma}[Domination by a reflected state]\label{lp:lem-fold-domination}
Fix \(\beta>0\) and an admissible coordinate or diagonal mirror
with normal \(n\) and height \(c\). For every bounded real Borel
function \(A\) of finitely many bonds and every
\(\nu\in\mathcal G_\beta^P\), there is
\(\rho\in\mathcal G_\beta^P\) with
\(\rho(A)\ge\nu(A)\). The state \(\rho\) is invariant under
\(\theta_c\) and reflection positive on continuous local functions
of bonds in the closed positive half-space.
\end{lemma}

\begin{proof}
Choose \(b\in P\) so that \(B=\tau_b A\) depends only on bonds strictly
above the mirror. Such a shift exists because \(n\) has a coordinate
equal to \(1\). Choose an even integer \(h>0\) large enough that
the endpoints of those bonds lie below height \(c+H\), where
\(w=hn\) and \(H=n\cdot w=h|n|^2\). Put
\[
 P_\parallel=P\cap n^\perp,\qquad
 Q=P_\parallel\oplus\mathbb Zw,\qquad
 Q_f=P_\parallel\oplus\mathbb Z(2w).
\]
These are subgroups of \(P\), with \(Q\) of finite index and
\([Q:Q_f]=2\). The tangential lattice \(P_\parallel\) has an
orthogonal integer basis. For a coordinate normal use the other
three doubled coordinate vectors. For \(n=e_\mu\pm e_\lambda\)
use \(2(e_\mu\mp e_\lambda)\) and the two doubled coordinate
vectors outside that pair.

Take the torus with period lattice
\(\Gamma_{M,J}=JP_\parallel\oplus\mathbb Z(2Mw)\), for sufficiently
large positive integers \(M,J\), and write \(\Gamma=\Gamma_{M,J}\).
Choose this cofinal tail so that the explicit diagonal period bounds
preceding Proposition~\ref{prop:lp-finite-rp} hold whenever needed.
It has \(2M\) slabs of height \(H\) between faces \(c+kH\).
All its side lengths tend to infinity as \(M,J\to\infty\).
For a local function \(D\) write \(D^\Gamma\) for its evaluation
on the periodic lift of a torus configuration. Consider the
unfolded and folded sums
\begin{align*}
 S_{\rm u}
   &=\sum_{v\in Q/\Gamma}(\tau_v B)^\Gamma,\\
 S_{\rm f}
   &=\sum_{v\in Q_f/\Gamma}
      \left((\tau_v B)^\Gamma+
            (\tau_v\theta_{c+H}B)^\Gamma\right).
\end{align*}
The first places translated copies in all slabs. The second
places translated copies in the even slabs and reflected copies
in the odd slabs. It is invariant under \(Q_f\) and all face
reflections. In fact translations by \(2w\) and reflection at
height \(c+H\) generate the normal action used to place its terms.

Let \(Z_{\rm u}(t)\) and \(Z_{\rm f}(t)\) be the Wilson partition
functions on this torus with added exponents \(tS_{\rm u}\) and
\(tS_{\rm f}\), respectively. To apply Lemma~\ref{lem:lp-chessboard},
use the slab function
\[
 f_k=\exp\!\left[t\sum_{v\in P_\parallel/JP_\parallel}
                          (\tau_{v+kw}B)^\Gamma\right]
\]
in slab \(k\). Its dissemination is a normal translate of the
folded weight. The unperturbed torus measure is invariant under
that translation, so each disseminated partition function is
\(Z_{\rm f}(t)\). The chessboard inequality gives
\begin{equation}\label{lp:eq-fold-partition}
                         Z_{\rm u}(t)\le Z_{\rm f}(t).
\end{equation}
The same finite comparison holds when \(B\) depends on bonds
of the closed slab, including bonds on its faces: the slab
chessboard lemma permits these shared variables and retains each
source occurrence with its own slab index.
For a test \(A\) straddling the mirror, the unfolded partition
function of its repeated array equals \(Z_{\rm u}(t)\) by the
global torus translation by \(b\).

Let \(\mu_{\Gamma,t}\) denote the normalized folded measure,
\(N_\Gamma=[\mathbb Z^4:\Gamma]\), and
\(\kappa=[\mathbb Z^4:Q]^{-1}\). Define
\[
 p_\Gamma(t)=N_\Gamma^{-1}
                \log\!\left(\frac{Z_{\rm f}(t)}{Z_{\rm f}(0)}\right).
\]
Finite-volume differentiation gives
\(p_\Gamma'(t)=\kappa\mu_{\Gamma,t}(B^\Gamma)\). Indeed the
\(\kappa N_\Gamma\) source terms all have the same mean by
\(Q_f\) periodicity and reflection at height \(c+H\).
Convexity and \(p_\Gamma(0)=0\) imply
\(p_\Gamma(t)\le t p_\Gamma'(t)\) for \(t>0\).
Apply Lemma~\ref{lem:lp-lower-tangent} to the unfolded array,
using \(\Gamma\subset Q\subset P\). For every fixed \(t>0\),
\begin{equation}\label{lp:eq-fold-pressure}
 t\kappa\nu(B)-o(1)
 \le N_\Gamma^{-1}\log\frac{Z_{\rm u}(t)}{Z_{\rm u}(0)}
 \le p_\Gamma(t)
 \le t\kappa\mu_{\Gamma,t}(B^\Gamma),
\end{equation}
where the error tends to zero as \(M,J\to\infty\).

At each pair of opposite slab faces the folded exponent has the
form \(E+\theta E\), with \(E\) its contribution from the slabs in
one closed half and \(\theta\) the reflection exchanging the halves.
Thus Proposition~\ref{prop:lp-finite-rp} gives reflection positivity
of \(\mu_{\Gamma,t}\) at these faces, by inserting \(Fe^{E}\) in
the unperturbed form. Its reflection invariance follows from the
same decomposition and invariance of the unperturbed measure.

At fixed \(0<t\le1\), apply Lemma~\ref{lem:lp-compactness} to
the bounded finite-range folded interaction. A subsequence of
\(\mu_{\Gamma,t}\) converges to a perturbed DLR state \(\rho_t\),
with convergence of bounded Borel local expectations.
The period \(Q_f\) and reflection invariance pass to the limit.
Every continuous local test in the closed half-space of height
at least \(c\) eventually lies in the corresponding torus half,
so its positivity passes to \(\rho_t\). From
\eqref{lp:eq-fold-pressure} we have \(\nu(B)\le\rho_t(B)\).

Now let \(t\downarrow0\) along a weakly convergent subsequence
of these states. On a fixed finite bond set, the added exponent
in the conditional density is \(O(t)\) uniformly in the exterior.
By Lemma~\ref{lem:lp-compactness}, the limit \(\rho\) is a Wilson
DLR state and \(\rho_t(B)\) converges to \(\rho(B)\) on this
subsequence. The state \(\rho\) retains \(Q_f\) periodicity and
the stated reflection properties, and \(\nu(B)\le\rho(B)\).
Theorem~\ref{lp:thm-link-symmetry} upgrades its period to \(P\).
Both \(\rho\) and \(\nu\) are then invariant under the original
shift \(b\), giving \(\nu(A)\le\rho(A)\). The volume limit was
taken at fixed \(t\) to remove the error in
\eqref{lp:eq-fold-pressure} before division by \(t\kappa\).
\end{proof}

\begin{samepage}
\begin{theorem}[All admissible mirrors]\label{lp:thm-all-mirrors}
At any fixed spacing \(a>0\) and coupling \(\beta\ge0\), every
Wilson \(\operatorname{SU}(2)\) DLR state with a finite-index
period subgroup is invariant under all lattice translations and
signed coordinate permutations. It is invariant and reflection
positive at each mirror
\[
 x_\mu=c,\quad c\in\tfrac a2\mathbb Z,\qquad
 x_\mu\pm x_\lambda=c,\quad c\in a\mathbb Z,\quad\mu<\lambda .
\]
The reflection inequality holds for every bounded Borel function
of the bonds whose endpoints lie in either closed half-space.
\end{theorem}
\end{samepage}

\begin{proof}
If \(\beta=0\), the DLR equation on each finite bond set forces
its marginal to be product Haar. Hence the state is the unique
product-Haar law. Permutation and inversion invariance of Haar
give all the stated lattice symmetries. Condition on the links
contained in the mirror. The remaining link sets of the two closed
halves are disjoint and independent, and the reflection fixes each
conditioned link. A link with exactly one endpoint on the mirror
belongs to the remaining set on its side. At a bond cut the crossing
links belong to neither half. Integrating them contributes one. The reflection form is therefore
an integral of an absolute square. This proves positivity for
continuous local tests, and Lemma~\ref{lem:lp-borel-rp} gives the
Borel assertion.

Assume \(\beta>0\). In lattice units put
\(K=\mathcal G_\beta^P\). For a fixed
mirror let \(K_c\) consist of the states in \(K\) invariant
under \(\theta_c\) and positive on the continuous local tests
of the closed positive half-space. The DLR and period equations
are closed affine conditions on probability measures. They can
be tested by continuous cylinder functions, which the Wilson
conditional kernels preserve. Hence \(K\) is compact and convex.
The reflection equations and positivity
inequalities on the specified tests give a closed convex
subset \(K_c\). It is nonempty by taking an unperturbed
subsequential limit on the split tori.

If \(\nu\in K\setminus K_c\), a weak neighborhood of \(\nu\)
defined by finitely many real continuous local tests is
disjoint from \(K_c\). Their joint moment vectors for
\(K_c\) form a compact convex subset of Euclidean space that
excludes the vector for \(\nu\). A strict separating linear
functional on these vectors gives a real continuous local
function \(A\) with \(\nu(A)>\sup_{\rho\in K_c}\rho(A)\).
This contradicts Lemma~\ref{lp:lem-fold-domination}, so \(K=K_c\).

Lemma~\ref{lem:lp-borel-rp} extends positivity to the bounded
Borel algebras of both closed half-spaces.

The mirror was arbitrary. Composing the site and bond reflections
at heights \(0\) and \(1/2\) in a coordinate direction gives a
unit translation. Site reflections through the origin give the
coordinate sign changes, and reflection in \(x_\mu=x_\lambda\)
gives the coordinate transpositions. These generate the stated
symmetries. The dilation \(x\mapsto ax\) restores the physical
spacing and the listed mirror heights.
\end{proof}

\section{Domination by finite chessboard seminorms}\label{lp:sec-board-domination}

For the cutoff boards of the model, let \(L\ge3\) be odd,
\(a=L^{-K}\), and let the fine board side be \(l=L^j\ge a\),
with \(K\) a nonnegative integer and \(j\) an integer. The
board faces cut lattice bonds in half. A cell function uses
only bonds whose endpoints lie in the open cell. On a cubic
torus of side \(\ell=2L^m\) with an integer \(m\ge\max(0,j)\), the
board has \(N=(\ell/l)^4\) cells and an even number in every
direction. Write \(p_{l,\ell}\) for the seminorm of
Corollary~\ref{cor:lp-board-chessboard} on this board.
Copies use conjugation with each orientation-reversing
isometry. A cell function may be presented on any actual
cell and transported to the reference cell to compute its
seminorm. The result is the same if the reference cell is
changed, by invariance of the torus measure and conjugation
invariance of the seminorm.

\begin{samepage}
\begin{theorem}[Fixed-spacing chessboard domination]\label{thm:lp-board-domination}
Fix \(\beta\ge0\), a periodic Wilson DLR state \(\rho\), and a
finite collection of occupied cells with bounded complex Borel
functions \(F_C\) on those cells. For every sequence of the
compatible torus sides \(\ell_k\to\infty\),
\begin{equation}\label{eq:lp-board-domination}
 \left|\rho\!\left(\prod_C F_C\right)\right|
    \le \liminf_{k\to\infty}\prod_C p_{l,\ell_k}(F_C).
\end{equation}
Only the occupied cells enter the products.
\end{theorem}
\end{samepage}

\begin{proof}
First let \(H\) be bounded, real and Borel on one cell. Enclose
a board of side \(\ell\) in a cube with faces on the board
faces, and sum the reflected copies of \(H\) over its cells.
Denote the sum by \(W_\ell\). Theorem~\ref{lp:thm-all-mirrors}
gives \(\rho(W_\ell)=N\rho(H)\).
Write \(\gamma\) for the DLR conditional kernel on the positive
bonds anchored at the lattice sites of the cube. Conditional
Jensen gives
\[
               N\rho(H)\le \rho\bigl(\log\gamma(e^{W_\ell})\bigr).
\]
Identify the anchored bonds with the torus bonds by reducing
their anchors modulo the period. The conditional interaction
and the torus action agree off a boundary layer. They differ
in at most \(C(\ell/a)^3\) plaquette terms, uniformly in the
outside configuration. The source copies use only bonds in
open cells, so their sum is the same function of the integrated
bonds in both integrals. Comparing the numerator and
denominator of the conditional exponential moment gives
\begin{equation}\label{eq:lp-board-jensen}
 \rho(H)\le \log p_{l,\ell}(e^H)
                   +C\beta\frac{l^4}{a^3\ell}.
\end{equation}
The error is divided by the number of cells. The comparison
uses bounded plaquette actions and holds at fixed \(a,\beta,l\).

For a bounded real cell function \(G\), take
\(H=\log(1+tG)\) with \(t\|G\|_\infty<1/2\) and \(t>0\).
Then \(\rho(H)\ge t\rho(G)-Ct^2\|G\|_\infty^2\).
The seminorm triangle inequality and \(p_{l,\ell}(1)=1\) give
\[
 \log p_{l,\ell}(1+tG)\le t p_{l,\ell}(G).
\]
Insert these in \eqref{eq:lp-board-jensen}. Take the lower
limit along the prescribed sequence at each fixed \(t\) and
then send \(t\) to zero. This proves
\(\rho(G)\le\liminf_k p_{l,\ell_k}(G)\).
For a complex cell function \(F\), choose a real phase
\(\varphi\) so that \(\rho(\operatorname{Re}(e^{-i\varphi}F))
=|\rho(F)|\). Homogeneity, the triangle inequality and
\(p(\overline F)=p(F)\) bound the seminorm of this real part
by \(p(F)\). We have proved \eqref{eq:lp-board-domination}
for a single cell.

For a finite collection of cells, choose an origin for the
coarse board on fine-board faces. Choose its allowed side
\(l'=L^r l\), with an integer \(r\ge1\) large enough that one coarse
cell contains all occupied fine cells strictly inside.
The coarse origin can be shifted independently of the fine
origin. The coarse faces are fine-board faces. The
coarse board is compatible with every sufficiently large
torus in the chosen sequence. The product \(G=\prod_C F_C\)
is one bounded complex function on the coarse cell.
Disseminate \(G\) on the coarse board. Every copied fine
factor is transported by an element of the fine-board
reflection group, since reflections in coarse faces belong
to that group. Its fine seminorm is therefore the seminorm
of the original factor. Each occupied-cell factor occurs
once in every coarse copy. The fine-board chessboard bound,
followed by the root for the number of coarse cells, yields
\[
                p_{l',\ell}(G)\le\prod_C p_{l,\ell}(F_C).
\]
Apply the one-cell bound to \(G\) on the coarse board and
take the lower limit to conclude.
\end{proof}

\begin{thebibliography}{Ba87}
\small
\bibitem[Ba87]{Bal87}
T.~Ba{\l}aban, \emph{Renormalization group approach to lattice gauge
field theories. I. Generation of effective actions in a small field
approximation and a coupling constant renormalization in four
dimensions}, Comm. Math. Phys. \textbf{109} (1987), 249--301.
\href{https://doi.org/10.1007/BF01215223}{doi:10.1007/BF01215223}.
\end{thebibliography}
\end{document}
